\scanpage{122}\mspage{77}
or, what comes to the same thing, the values
\[
\handwritten{\frac{\unsure{\Phi|\phi}_{\unsure{0|o}}}{\Phi_k}(Q),\ldots\,.,\frac{\Phi_T}{\Phi_k}(Q)}
\]
from which it follows that
\[
\handwritten{\bigl(\Phi_0^{(i)}(Q'),\ldots\,.,\Phi_T^{(i)}(Q')\bigr)}=\handwritten{\bigl(x_0(Q),\ldots\,.,x_T(Q)\bigr)}
\]
\margin{p. 210}and
\[
\handwritten{Q}=\handwritten{f_nQ'}
\]

It is easy to show that $f_n$ is a continuous mapping
(in Zariski's topology) and even a regular one. The latter
follows from the fact that, in $Q\in V^{(n)}$ ($x_k\ne 0$ at $Q$),
any regular function $R$ has the form
\[
\handwritten{R}=\handwritten{\frac{F_1\left(\frac{x_0}{x_k},\ldots,\frac{x_T}{x_k}\right)}{F_2\left(\frac{x_0}{x_k},\ldots\,.\frac{x_T}{x_k}\right)}}
\]
where $F_1$ and $F_2$ are polynomials and $F_2(Q)\ne0$; the
corresponding function on $V$ will be
\[
\handwritten{R'}=\handwritten{\frac{F_1\left(\frac{\Phi_0}{\Phi_k},\ldots\,.,\frac{\Phi_T}{\Phi_k}\right)}{F_2\left(\frac{\Phi_0}{\Phi_k},\ldots\,.,\frac{\Phi_T}{\Phi_k}\right)}}
\]
% archive sheet number: 121
